\chapter{Inspection and recording}
\label{ch:t-inspection}

\section{At a stop point}

\code{inspection_node} runs when the follower reaches a stop point: it finds the inspection
points of the site within \param{match_radius_m} (\SI{1.5}{m}) of the stop, and for each target
moves the PTZ camera, waits, measures and judges.

\begin{center}
\includegraphics[width=\linewidth]{inspection}
\end{center}

\begin{center}
\begin{tabularx}{\linewidth}{@{}>{\raggedright}p{0.32\linewidth}p{0.12\linewidth}X@{}}
  \toprule
  \textbf{Setting} & \textbf{Default} & \textbf{Tune when} \\
  \midrule
  \param{camera.move_timeout_sec} & \SI{20}{s} & a long pan fails: the camera is slow in the cold \\
  \param{camera.settle_sec} & \SI{1}{s} & images blur or the focus hunts: raise to \SIrange{2}{3}{s} at high zoom \\
  \param{thermal.samples} & 3 & readings jump: more samples (the median is used) \\
  \param{thermal.sample_interval_sec} & \SI{0.3}{s} & the frames are the same: longer \\
  \param{max_duration_sec} & \SI{150}{s} & a point has many targets: later targets are skipped \\
  \param{acoustic.duration_sec} & \SI{5}{s} & the acoustic camera needs longer to classify \\
  \bottomrule
\end{tabularx}
\end{center}

\begin{notebox}
\param{max_duration_sec} must stay below the follower's \param{hooks.inspection_timeout_sec}
(\SI{180}{s} on the X30): otherwise the robot walks on while the camera still works.
\end{notebox}

\section{Thermal criteria}

Each target has one or more regions on the thermal image. With $T_{\max}$ the hottest region and
$T_a$ the air temperature (\param{thermal.ambient_c}), three differences are judged against their
levels, after NETA MTS practice:
\begin{align}
  \text{absolute:}&\quad T_{\max}\\
  \text{ambient:}&\quad \Delta T_a = T_{\max} - T_a\\
  \text{similar:}&\quad \Delta T_s = T_{\max} - \operatorname{median}\{T_j : j \ne \max\}
    \qquad\text{(two regions or more)}
\end{align}

\begin{center}
\begin{small}
\begin{tabular}{@{}lrrr@{}}
  \toprule
  \textbf{Criterion} & \textbf{Attention} & \textbf{Warning} & \textbf{Critical} \\
  \midrule
  absolute $T_{\max}$ & \SI{80}{\celsius} & \SI{90}{\celsius} & \SI{105}{\celsius} (or the target's \param{max_temp_c}) \\
  ambient $\Delta T_a$ & \SI{11}{K} & \SI{21}{K} & \SI{41}{K} \\
  similar $\Delta T_s$ & \SI{2}{K} & \SI{4}{K} & \SI{16}{K} \\
  \bottomrule
\end{tabular}
\end{small}
\end{center}

\noindent The target's severity is the highest of its three. A target taught with a rated
temperature (\param{max_temp_c} in NETRA) uses it as its absolute critical level.

\begin{center}
\includegraphics[width=\linewidth]{severity}
\end{center}

\begin{warnbox}[The air temperature is live: set it]
$\Delta T_a$ is only as good as $T_a$. The default \SI{30}{\celsius} misjudges a hot afternoon or a
cool night by the difference. Set it each day (it takes effect within a second):
\cmd{rosparam set /inspection_node/thermal/ambient_c 34}.
\end{warnbox}

\subsection{Getting good temperatures}

\begin{itemize}
  \item \textbf{Emissivity}: the camera's setting applies to every target. Bare, oxidized copper and
    aluminium read far too cold; the similar criterion compares like with like and suffers least.
  \item \textbf{Distance and zoom}: a region needs several pixels of the hot spot; zoom in from the
    stop rather than teach a closer stop next to the equipment.
  \item \textbf{Sun and reflections}: a reflection of the sun on a busbar is hot to the camera. Patrol
    thermal routes at night or on a cloudy day where it matters.
  \item \textbf{Load}: a joint's heating rises with the square of the current. Readings at light
    load understate a defect; NETRA's trend over patrols shows it best.
\end{itemize}

\section{Partial discharge}

The acoustic camera classifies what it hears; each class maps to a severity:
\param{acoustic.corona} (attention), \param{acoustic.surface} (warning), \param{acoustic.floating}
(critical), noise is normal. \param{intensity_warning} and \param{intensity_critical} raise a
discharge by its intensity as well (0: off): set them after a few patrols, from the intensities
the site's normal corona shows.

\section{Reports}

Every point inspected writes a folder under \param{reports_dir} with the readings, the images and,
with \param{camera.raw_thermal}, the raw temperatures per pixel. The oldest days go while less
than \param{reports_min_free_gb} (\SI{6}{GB}) is free; \param{reports_keep_days} bounds them by age.

\section{The recorder}

\code{recorder_node} (perception host) records the cameras during every patrol (\param{patrol.enabled}) and on
request: the visible and thermal streams as MP4 segments of \param{video.segment_sec}, every raw
thermal frame with the robot's pose and the camera's pan, tilt and zoom, and the acoustic data in
chunks.

\begin{center}
\begin{tabularx}{\linewidth}{@{}>{\raggedright}p{0.3\linewidth}p{0.12\linewidth}X@{}}
  \toprule
  \textbf{Setting} & \textbf{Default} & \textbf{Meaning} \\
  \midrule
  \param{disk.min_free_gb} & \SI{10}{GB} & the oldest sessions go while less is free \\
  \param{disk.stop_free_gb} & \SI{3}{GB} & below this nothing more is written \\
  \param{disk.check_sec} & \SI{30}{s} & checked this often while recording \\
  \param{patrol.stop_after_sec} & \SI{30}{s} & after the route ends; laps stay one session \\
  \param{thermal.min_interval_sec} & 0 & every thermal frame; raise it to save disk \\
  \bottomrule
\end{tabularx}
\end{center}

\noindent A raw thermal frame holds a temperature for every pixel: at the camera's full frame
rate an hour of patrol fills gigabytes. Set \param{thermal.min_interval_sec} to \SI{1}{s} or more
when the raw frames serve trends rather than video.
